%=============================================================================
\section{Results}
\label{sec:results}
%=============================================================================

Results are reported in the order of the validation design
(Section~\ref{sec:step5}): what the generator reproduces, how it compares with
an independent survey, and then, in Section~\ref{sec:scenarios}, the projection
for the outcome the national survey does not measure. All figures can be
recomputed with the scripts of Appendix~\ref{app:code}.

%-----------------------------------------------------------------------------
\subsection{Reproduction of observed variables}
\label{sec:res_fidelity}
%-----------------------------------------------------------------------------

Across 78 alignable variables, the synthetic population has a mean total
variation distance (TVD) of %%TVDMEAN%% and a median of %%TVDMEDIAN%% from the
CIS marginals; %%TVD05%% of variables fall below $0.05$ and %%TVD10%% below
$0.10$. Table~\ref{tab:comparison} compares BVS with three alternatives fitted
to the same data.

\begin{table}[htbp]
\centering
\caption{Synthetic-data methods on CIS~3391 (78 alignable variables). For
observed variables BVS is a Bayesian network with a BDeu prior and
survey-weighted counts, as in \citet{martins2024}. $\chi^2$: dependencies significant in the CIS that remain
significant in the synthetic data, of those tested. $V$ ratio: range of
$V_{\mathrm{synthetic}} / V_{\mathrm{CIS}}$ over the tested pairs; 1 indicates
preserved strength. CTGAN, TVAE and the independent baseline are from the
deposited comparison (v1.0).}
\label{tab:comparison}
\small
\begin{threeparttable}
\begin{tabular}{lccccccc}
\toprule
\textbf{Method} & \textbf{TVD}$_\mu$ & \textbf{TVD}$_{\tilde x}$ &
  \textbf{TVD${<}0.05$} & \textbf{Corr.} & $\chi^2$ & \textbf{$V$ ratio} &
  \textbf{Unobserved} \\
 & $\downarrow$ & $\downarrow$ & \textbf{(\%)} $\uparrow$ &
  \textbf{MAE} $\downarrow$ & \textbf{kept} & \textbf{(range)} &
  \textbf{outcomes} \\
\midrule
%%CMPROWS%%
\bottomrule
\end{tabular}
\begin{tablenotes}\small
\item The independent baseline samples each marginal separately; it reproduces
  marginals by construction and no dependency. A random-forest detector
  separates every synthetic set from the CIS (AUC $\geq 0.99$ for all
  methods); the detector was trained on unbalanced classes and the figure is
  not informative about relative quality.
\end{tablenotes}
\end{threeparttable}
\end{table}

BVS has a mean TVD 4.3 and 6.0 times lower than CTGAN and TVAE, and a median
TVD 11 and 12 times lower. CTGAN shrinks association strength to 10--31\% of its CIS
value, and TVAE inflates it by up to a factor of six.

Table~\ref{tab:vratios} reports dependency preservation pair by pair.
Dependencies within thematic blocks are kept (environmental concern with
recycling, %%VREC%%), whereas those that cross blocks, in particular between
demographics and attitudes, are attenuated to between a fifth and a half of
their CIS strength; %%KEPT%% of the ten remain statistically detectable. This
follows from block-wise structure learning: variables in different blocks are
linked only through the cross-block screening of Step~1. Joint distributions
across blocks are therefore less faithful than marginals. For the target
outcome this matters less, because the transfer model is evaluated on each
synthetic individual's own characteristics; it matters for any analysis that
relies on associations among observed attitudes. %%SEEDTEXT%%

\begin{table}[htbp]
\centering
\caption{Cram\'er's $V$ in the CIS and in the BVS synthetic population for the
ten tested dependencies. Ratio $= V_{\mathrm{BVS}} / V_{\mathrm{CIS}}$.
\emph{Kept}: the dependency remains significant at the 5\% level.}
\label{tab:vratios}
\small
\begin{tabular}{lcccc}
\toprule
\textbf{Dependency} & $V_{\mathrm{CIS}}$ & $V_{\mathrm{BVS}}$ &
  \textbf{Ratio} & \textbf{Kept} \\
\midrule
%%VROWS%%
\bottomrule
\end{tabular}
\end{table}


%-----------------------------------------------------------------------------
\subsection{External comparison with an independent Spanish survey}
\label{sec:res_ipsos}
%-----------------------------------------------------------------------------

One item of the Ipsos survey for Spain \citep{ipsos2022} is loosely comparable
with a CIS variable: support for an international treaty on plastics against
environmental concern (V15, top three of five categories). The synthetic share
is %%IPSOSBVS%% against 90\% in Ipsos, a difference of %%IPSOSERR%%; the CIS
itself gives %%IPSOSCIS%%. The difference therefore reflects the two
instruments, not the generator. The other Ipsos items have no counterpart in
the CIS: an earlier comparison of manufacturer responsibility with V27 matched
an item on paying higher taxes, coded from ``strongly in favour'' to ``strongly
against'', and is discarded. This comparison is a plausibility check under
loose construct matching and does not bear on food-packaging acceptance.

